% !TeX program = xelatex
% Самодостатній файл UTF-8. Компіляція: xelatex <ім’я-файлу>.tex (двічі).
\documentclass[10pt,a4paper,landscape]{article}
\usepackage[margin=15mm]{geometry}
\usepackage{fontspec}
\setmainfont{FreeSerif.otf}[BoldFont=FreeSerifBold.otf,ItalicFont=FreeSerifItalic.otf,BoldItalicFont=FreeSerifBoldItalic.otf]
\usepackage{polyglossia}
\setdefaultlanguage{ukrainian}
\usepackage{longtable,array,booktabs,xcolor,fancyhdr}
\usepackage[hidelinks]{hyperref}
\definecolor{accent}{HTML}{16656B}
\pagestyle{fancy}\fancyhf{}
\fancyhead[L]{STEM • 5 клас}\fancyhead[R]{Навчальний та календарно-тематичний план}
\fancyfoot[C]{\thepage}
\setlength{\headheight}{14pt}
\setlength{\parindent}{0pt}
\setlength{\parskip}{5pt}
\renewcommand{\arraystretch}{1.18}
\newcolumntype{P}[1]{>{\raggedright\arraybackslash}p{#1}}
\begin{document}
{\LARGE\color{accent} Навчальний план STEM, 5 клас}

{\large 1,5 год/тиждень • 48 годин на рік}

Заклад освіти: \rule{90mm}{0.3pt}\quad Навчальний рік: \rule{30mm}{0.3pt}

Учитель/учителька: \rule{95mm}{0.3pt}\quad Клас: \rule{25mm}{0.3pt}
\section*{Пояснювальна записка}
Авторський адаптований план укладено за наданою модельною програмою «STEM. 5--6 класи (міжгалузевий інтегрований курс)» О. В. Бутурліної та О. Є. Артєм’євої (розділ 5 класу, с. 9--20 PDF) і зошитом-конспектом «STEM-LAB. 5» (О. Бутурліна та ін., за загальною редакцією О. В. Коршунової, ВД «Освіта», 2023; файл «Stem. Full.pdf»).
Мета: формувати дослідницькі, інженерні, математичні й цифрові вміння, навички співпраці та уявлення про STEM-професії через практичні проєкти.

У наданій програмі орієнтовний розподіл становить 35 годин: 3 + 7 + 5 + 5 + 5 + 5 + 5. Програма дозволяє вчителеві змінювати розподіл годин зі збереженням очікуваних результатів (с. 8 PDF). Цей план адаптовано до 32 навчальних тижнів. Збережено вступ, п’ять змістових модулів і підсумкову STEM-практику. Послідовність модулів відповідає програмі; у зошиті розділ про сили передує розділу про Всесвіт.
\textbf{Облік часу:} план містить 48 окремих уроків по 1 навчальній годині. Протягом двотижневого циклу проводять 3 уроки: один урок першого тижня та два уроки другого. І семестр (15 тижнів) — 22 уроки; ІІ семестр (17 тижнів) — 26 уроків; разом 48 уроків. Середнє навантаження становить 1,5 години на тиждень. Кожен рядок календарної таблиці відповідає одному уроку з власною назвою; дати вчитель/учителька вписує відповідно до розкладу. Практикуми об’єднують практичні частини двох сусідніх тематичних блоків у повний урок.

\textbf{Ресурси та безпека:} зошит, папір, картон, лінійки, фарби, ліхтарик, низьковольтні навчальні набори; за можливості — комп’ютер зі Scratch і камера. Цифрові завдання мають паперовий варіант. Електричні моделі працюють лише від батарейок; не використовувати мережеву напругу. Не спрямовувати оптичні прилади на Сонце або очі. Запуски моделей проводити у вільній зоні. Складні вирізи готує вчитель.

\textbf{Джерельні позначення:} «Зошит» у таблиці — номер сторінки наданого 64-сторінкового скану, з обкладинкою як сторінкою 1. Це опора для добору вправ, а не вимога виконати всі завдання сторінок. Частину практикумів, Scratch, семестрову діагностику й хакатон адаптовано або запропоновано за програмою; вони не подаються як дослівні завдання зошита.
\newpage
\section*{Розподіл навчального часу}
\begin{tabular}{P{125mm}rrr}
\toprule Розділ & І семестр & ІІ семестр & Разом \\\midrule
Вступ & 4 & 0 & 4 \\
Людина — людина. Я у школі & 9 & 0 & 9 \\
Людина — природа. Я у Всесвіті & 9 & 0 & 9 \\
Людина — техніка. Сила — це сила & 0 & 8 & 8 \\
Людина — образ. Промінь і світло & 0 & 6 & 6 \\
Людина — знак. Під знаком STEM & 0 & 6 & 6 \\
Хакатон, фестиваль, STEM-практика & 0 & 6 & 6 \\
\midrule Разом & 22 & 26 & 48 \\
\bottomrule\end{tabular}
\section*{Оцінювання та очікувані результати}
На першому занятті — стартова бесіда й короткі практичні завдання без окремої додаткової години. Упродовж навчання — спостереження за роботою, аналіз таблиць, ескізів, моделей, короткі пояснення учнів. Після кожного проєкту — самооцінювання та взаємний відгук. Семестрові підсумки інтегровані у тижні 15 і 32.

Спільні критерії: зрозуміла мета; коректні спостереження й розрахунки; працездатність та безпечність моделі; пояснення змін після випробування; співпраця й особистий внесок. Для формувального відгуку використовуються описи «виконує з допомогою», «виконує самостійно», «пояснює та переносить у нову ситуацію». Конкретну шкалу підсумкового оцінювання визначає заклад освіти.

Наприкінці року учень/учениця формулює просту проблему й гіпотезу, планує послідовність дій, вимірює та подає дані в таблиці або діаграмі, створює і вдосконалює модель, користується знаками та цифровими засобами, пояснює рішення і називає пов’язані професії. Портфоліо містить особистий план, результати п’яти тематичних проєктів і підсумковий прототип.
\newpage
\section*{Календарно-тематичний план. І семестр}
\fontsize{8.2}{9.6}\selectfont\begin{longtable}{P{10mm}P{11mm}P{14mm}P{42mm}P{17mm}P{78mm}P{46mm}}
\toprule Урок & Тиждень & Дата & Модуль і тема & Зошит & Навчальна діяльність & Очікуваний результат \\\midrule\endfirsthead
\toprule Урок & Тиждень & Дата & Модуль і тема & Зошит & Навчальна діяльність & Очікуваний результат \\\midrule\endhead
\bottomrule\endfoot
1 & 1 & & \textbf{Вступ}\par STEM-освіта та світ професій & 2--5 & Коло асоціацій STEM; сортування професій; стартова бесіда. & Пояснює складники STEM і наводить приклади професій. \\ \midrule
2 & 2 & & \textbf{Вступ}\par Проєкт: від задуму до плану & 6--9 & Сформулювати мету шкільного STEM-центру, розподілити ролі, скласти резюме проєкту. & Визначає мету, продукт, ресурси й етапи проєкту. \\ \midrule
3 & 2 & & \textbf{Вступ}\par Практикум: діаграма STEM-професій і вибір проєктної ідеї & 2--5; 6--9 & Скласти діаграму професій за результатами опитування класу. \par Порівняти два задуми STEM-центру за доступними ресурсами. & Застосовує набуті вміння у практичній роботі; фіксує результат і пояснює висновок. \\ \midrule
4 & 3 & & \textbf{Вступ}\par Мій шлях до успіху & 6--7 & Створити особистий план розвитку; обговорити рівні можливості у виборі професії. & Визначає власну навчальну ціль і спосіб перевірки поступу. \\ \midrule
5 & 4 & & \textbf{Людина — людина. Я у школі}\par Моя школа: люди та професії & 10--12 & Мініопитування про шкільні професії; таблиця й діаграма складу класу. & Збирає знеособлені дані й пояснює діаграму. \\ \midrule
6 & 4 & & \textbf{Людина — людина. Я у школі}\par Практикум: карта сильних сторін та інтерв’ю зі шкільним фахівцем & 6--7; 10--12 & Додати до портфоліо карту сильних сторін і взаємний відгук. \par Підготувати запитання для інтерв’ю зі шкільним працівником. & Застосовує набуті вміння у практичній роботі; фіксує результат і пояснює висновок. \\ \midrule
7 & 5 & & \textbf{Людина — людина. Я у школі}\par Територія школи. Масштаб і план & 13--14 & Виміряти клас; обчислити площу; виконати план у вибраному масштабі. & Переносить виміри на план і застосовує умовні позначення. \\ \midrule
8 & 6 & & \textbf{Людина — людина. Я у школі}\par Умови та засоби навчання & 15--16 & Порівняти навчальні засоби різних часів; спланувати дослідження класу. & Обґрунтовує вибір засобів навчання за заданими критеріями. \\ \midrule
9 & 6 & & \textbf{Людина — людина. Я у школі}\par Практикум: точність плану школи й дослідження освітленості & 13--14; 15--16 & Порівняти план із картою; знайти та виправити похибки вимірювання. \par Провести командне дослідження освітленості на якісному рівні. & Застосовує набуті вміння у практичній роботі; фіксує результат і пояснює висновок. \\ \midrule
10 & 7 & & \textbf{Людина — людина. Я у школі}\par Чому енергоефективність важлива & 17--19 & Порівняти споживання ламп за наданими даними; запропонувати способи економії. & Обчислює витрати енергії за зразком і пояснює вибір рішення. \\ \midrule
11 & 8 & & \textbf{Людина — людина. Я у школі}\par Енергоефективна сигнальна система & 20--21 & Скласти низьковольтне коло з батарейки, вимикача та світлодіодного модуля. & Читає просту схему, збирає і перевіряє модель за інструкцією. \\ \midrule
12 & 8 & & \textbf{Людина — людина. Я у школі}\par Практикум: діаграма енергоспоживання та випробування сигналізації & 17--19; 20--21 & Перевірити розрахунки для іншої тривалості освітлення; побудувати діаграму. \par Випробувати модель; знайти причину несправності та поліпшити з’єднання. & Застосовує набуті вміння у практичній роботі; фіксує результат і пояснює висновок. \\ \midrule
13 & 9 & & \textbf{Людина — людина. Я у школі}\par Школа моєї мрії: захист проєкту & 22--23 & Зібрати простий макет із підготовлених деталей; представити енергоощадне рішення. & Пов’язує план, модель і потреби користувачів; оцінює внесок команди. \\ \midrule
14 & 10 & & \textbf{Людина — природа. Я у Всесвіті}\par Об’єкти Всесвіту. Сонячна система & 36--38 & Побудувати порівняльну модель планет; обговорити масштаб і межі моделі. & Розрізняє основні об’єкти й пояснює умовність їх зображення. \\ \midrule
15 & 10 & & \textbf{Людина — природа. Я у Всесвіті}\par Практикум: паспорт шкільного проєкту й анімація руху у Всесвіті & 22--23; 36--38 & Доопрацювати макет за відгуками; оформити паспорт проєкту. \par Створити просту анімацію руху об’єкта в Scratch; за відсутності техніки — алгоритм на папері. & Застосовує набуті вміння у практичній роботі; фіксує результат і пояснює висновок. \\ \midrule
16 & 11 & & \textbf{Людина — природа. Я у Всесвіті}\par Зоряне небо та оптичні прилади & 39--40 & Виготовити зоряний ліхтарик за шаблоном; порівняти телескоп і бінокль. & Пояснює призначення оптичних приладів і демонструє модель. \\ \midrule
17 & 12 & & \textbf{Людина — природа. Я у Всесвіті}\par Від повітряної кулі до космічної станції & 41--42 & Створити стрічку освоєння космосу; змоделювати станцію з готових елементів. & Називає функції станції та обов’язки екіпажу. \\ \midrule
18 & 12 & & \textbf{Людина — природа. Я у Всесвіті}\par Практикум: оптичні спостереження та система життєзабезпечення станції & 39--40; 41--42 & Змінити відстань до екрана; записати, як змінюється зображення. \par Розробити схему життєзабезпечення станції й розподіл роботи екіпажу. & Застосовує набуті вміння у практичній роботі; фіксує результат і пояснює висновок. \\ \midrule
19 & 13 & & \textbf{Людина — природа. Я у Всесвіті}\par Людина у відкритому космосі & 43--44 & Порівняти захисні матеріали; спроєктувати модель захисту від холоду та ударів. & Пов’язує небезпеку з властивістю захисного матеріалу. \\ \midrule
20 & 14 & & \textbf{Людина — природа. Я у Всесвіті}\par Космічні професії. Україна космічна & 45--46 & Підготувати картку професії й інформаційний буклет за матеріалами вчителя. & Називає компетентності фахівця; зазначає джерела інформації. \\ \midrule
21 & 14 & & \textbf{Людина — природа. Я у Всесвіті}\par Практикум: захист у космосі та перевірка джерел інформації & 43--44; 45--46 & Порівняти два зразки захисту в безпечному модельному випробуванні. \par Перевірити твердження буклета за другим джерелом; додати ілюстративну схему. & Застосовує набуті вміння у практичній роботі; фіксує результат і пояснює висновок. \\ \midrule
22 & 15 & & \textbf{Людина — природа. Я у Всесвіті}\par Я у Всесвіті: представлення результатів & 36--46 & Представити космічні моделі та буклети; виконати коротку семестрову діагностику. & Пояснює рішення з опорою на дані; визначає, що варто повторити. \\ \midrule
\end{longtable}\normalsize
\newpage
\section*{Календарно-тематичний план. ІІ семестр}
\fontsize{8.2}{9.6}\selectfont\begin{longtable}{P{10mm}P{11mm}P{14mm}P{42mm}P{17mm}P{78mm}P{46mm}}
\toprule Урок & Тиждень & Дата & Модуль і тема & Зошит & Навчальна діяльність & Очікуваний результат \\\midrule\endfirsthead
\toprule Урок & Тиждень & Дата & Модуль і тема & Зошит & Навчальна діяльність & Очікуваний результат \\\midrule\endhead
\bottomrule\endfoot
23 & 16 & & \textbf{Людина — техніка. Сила — це сила}\par Сили у природі. Історія повітроплавання & 24--25 & Дослідити дію сили на рух і форму предмета; порівняти літальні апарати. & Наводить приклади дії сил і розрізняє спостереження та припущення. \\ \midrule
24 & 16 & & \textbf{Людина — техніка. Сила — це сила}\par Практикум: вдосконалення космічної моделі та схема дії сил & 36--46; 24--25 & Виправити помилки діагностики; удосконалити модель за відгуками. \par Скласти схему сил, що діють на модель літального апарата. & Застосовує набуті вміння у практичній роботі; фіксує результат і пояснює висновок. \\ \midrule
25 & 17 & & \textbf{Людина — техніка. Сила — це сила}\par Сила тертя & 26--28 & Порівняти рух однакового тіла на різних поверхнях; заповнити таблицю. & Робить висновок про тертя, змінюючи одну умову досліду. \\ \midrule
26 & 18 & & \textbf{Людина — техніка. Сила — це сила}\par Прості машини та механізми & 29--31 & Дослідити важіль на лінійці; знайти прості механізми у побутових предметах. & Пояснює дію важеля на основі досліду. \\ \midrule
27 & 18 & & \textbf{Людина — техніка. Сила — це сила}\par Практикум: дослідження тертя і робота важеля & 26--28; 29--31 & Повторити вимірювання; пояснити розбіжності та корисне застосування тертя. \par Порівняти положення опори й вантажу; оформити схему спостережень. & Застосовує набуті вміння у практичній роботі; фіксує результат і пояснює висновок. \\ \midrule
28 & 19 & & \textbf{Людина — техніка. Сила — це сила}\par Крила: конструювання моделі & 32--34 & Виготовити паперовий літак; висунути припущення про вплив форми крил. & Створює модель за ескізом і формулює перевірювану гіпотезу. \\ \midrule
29 & 20 & & \textbf{Людина — техніка. Сила — це сила}\par Крила: випробування й захист & 32--35 & Виміряти дальність повторних запусків; порівняти результати та представити висновок. & Фіксує результати й аргументує поліпшення моделі. \\ \midrule
30 & 20 & & \textbf{Людина — техніка. Сила — це сила}\par Практикум: варіанти крил і діаграма результатів запусків & 32--34; 32--35 & Створити другий варіант крил, зберігаючи інші умови. \par Побудувати діаграму запусків; повторно перевірити вдосконалену модель. & Застосовує набуті вміння у практичній роботі; фіксує результат і пояснює висновок. \\ \midrule
31 & 21 & & \textbf{Людина — образ. Промінь і світло}\par Світло і колір & 47--49 & Дослідити спектр на готовій моделі спектроскопа; порівняти кольорові зображення. & Описує спектральний склад білого світла й результати спостереження. \\ \midrule
32 & 22 & & \textbf{Людина — образ. Промінь і світло}\par Фарби та відтінки & 50--51 & Змішати фарби; створити палітру; порівняти природні барвники. & Отримує відтінки й відрізняє змішування фарб від змішування світла. \\ \midrule
33 & 22 & & \textbf{Людина — образ. Промінь і світло}\par Практикум: простий спектроскоп і рецепт кольору & 47--49; 50--51 & Зібрати простий спектроскоп за підготовленим шаблоном і порівняти безпечні джерела світла. \par Дослідити вплив співвідношення фарб; записати рецепт відтінку. & Застосовує набуті вміння у практичній роботі; фіксує результат і пояснює висновок. \\ \midrule
34 & 23 & & \textbf{Людина — образ. Промінь і світло}\par Поезія та фотографія: образ Сонця & 52--55 & Зіставити художній образ із явищем; створити фотографію або паперову композицію світла і тіні. & Добирає засоби зображення й пояснює роль освітлення. \\ \midrule
35 & 24 & & \textbf{Людина — образ. Промінь і світло}\par Проєкт «Намалюю тобі Сонце» & 47--55 & Оформити творчу роботу з палітрою, світлиною та поясненням світлового явища. & Поєднує наукове пояснення і художній задум; презентує результат. \\ \midrule
36 & 24 & & \textbf{Людина — образ. Промінь і світло}\par Практикум: фотопара та мініекспозиція «Намалюю тобі Сонце» & 52--55; 47--55 & Порівняти два варіанти освітлення одного предмета; оформити фотопару. \par Провести мініекспозицію та вдосконалити роботу за відгуками. & Застосовує набуті вміння у практичній роботі; фіксує результат і пояснює висновок. \\ \midrule
37 & 25 & & \textbf{Людина — знак. Під знаком STEM}\par Знаки й професії системи «людина — знак» & 56--58 & Класифікувати знаки; закодувати коротке повідомлення; обговорити професії. & Розрізняє природні й штучні знаки та пояснює правила кодування. \\ \midrule
38 & 26 & & \textbf{Людина — знак. Під знаком STEM}\par Графічні знаки. Геральдика & 57--60 & Розглянути герби; створити ескіз сімейного герба з поясненням символів. & Пояснює зв’язок символу зі змістом і читає графічні повідомлення. \\ \midrule
39 & 26 & & \textbf{Людина — знак. Під знаком STEM}\par Практикум: перевірка шифру та створення печатки & 56--58; 57--60 & Обмінятися шифрами; перевірити однозначність декодування. \par Виготовити просту печатку з безпечних готових матеріалів; перевірити відбиток. & Застосовує набуті вміння у практичній роботі; фіксує результат і пояснює висновок. \\ \midrule
40 & 27 & & \textbf{Людина — знак. Під знаком STEM}\par Кодування імені. STEM-символіка & 58--61 & Закодувати ім’я за обраною абеткою; створити ескіз емблеми. & Користується узгодженим кодом і обґрунтовує символи. \\ \midrule
41 & 28 & & \textbf{Людина — знак. Під знаком STEM}\par Проєкт «Я обираю STEM» & 61--63 & Оформити емблему та презентувати її значення; провести взаємооцінювання. & Передає ідеї STEM символами й пояснює власне рішення. \\ \midrule
42 & 28 & & \textbf{Людина — знак. Під знаком STEM}\par Практикум: варіанти STEM-емблеми й перевірка читабельності & 58--61; 61--63 & Перевірити код у парі; створити другий варіант емблеми. \par Перевірити читабельність емблеми різного розміру; завершити портфоліо. & Застосовує набуті вміння у практичній роботі; фіксує результат і пояснює висновок. \\ \midrule
43 & 29 & & \textbf{Хакатон, фестиваль, STEM-практика}\par STEM-практика: професії громади & 64 & Провести очну або віртуальну зустріч із фахівцем; за відсутності — дослідити підготовлений кейс. & Ставить змістовні запитання й пов’язує професію з навчальними вміннями. \\ \midrule
44 & 30 & & \textbf{Хакатон, фестиваль, STEM-практика}\par Хакатон: задум і прототип & 62--64 & Обрати проблему класу; спланувати й створити простий прототип із доступних матеріалів. & Визначає проблему, критерії успіху, ресурси й ролі. \\ \midrule
45 & 30 & & \textbf{Хакатон, фестиваль, STEM-практика}\par Практикум: картка STEM-організації та перше випробування прототипу & 64; 62--64 & Підготувати інформаційну картку місцевої STEM-організації. \par Провести першу перевірку прототипу за критеріями. & Застосовує набуті вміння у практичній роботі; фіксує результат і пояснює висновок. \\ \midrule
46 & 31 & & \textbf{Хакатон, фестиваль, STEM-практика}\par Хакатон: випробування і вдосконалення & 62--64 & Випробувати прототип; зафіксувати зміни; підготувати короткий виступ. & Обґрунтовує зміни результатами випробування. \\ \midrule
47 & 32 & & \textbf{Хакатон, фестиваль, STEM-практика}\par STEM-фестиваль. Підсумкова рефлексія & 62--64 & Представити портфоліо й прототип; виконати підсумкові завдання та самооцінювання. & Демонструє поступ у дослідженні, конструюванні та співпраці. \\ \midrule
48 & 32 & & \textbf{Хакатон, фестиваль, STEM-практика}\par Практикум: повторне випробування та особисті навчальні цілі & 62--64 & Повторити випробування та порівняти початковий і кінцевий результати. \par Обговорити індивідуальні результати й сформулювати цілі на наступний рік. & Застосовує набуті вміння у практичній роботі; фіксує результат і пояснює висновок. \\ \midrule
\end{longtable}\normalsize
\end{document}
